本文围绕神经网络特征学习的动力学特性展开研究，主要创新点如下：

\begin{enumerate}

    \item \textbf{构建了跨任务场景的统一特征学习分析范式}: 
    本文结合传统学习理论与特征学习理论的最新研究成果, 系统性地从通用任务 (对抗训练), 图像任务 (异常检测) 和文本任务 (问答系统) 三个维度探索神经网络的特征学习机制, 揭示了不同任务场景下特征提取, 演化与网络适配的共性规律. 
    通过建立统一的理论分析框架, 本文为理解 ``神经网络如何自动提取有效特征'', ``特征如何随训练过程演化'' 以及 ``为何特定结构的网络在特定任务中表现更优'' 这三个核心问题提供了系统性的理论回答, 推动了特征学习理论在不同应用场景中的统一与发展。
    
    \item \textbf{发现并理论阐释了对抗训练下的单纯形压缩现象}: 
    本文首次系统性地揭示了对抗训练 (Adversarial Training, AT) 对神经网络特征几何结构的影响, 发现了单纯形压缩 (Simplex Compression) 这一新现象, 并基于特征学习分析框架证明了该现象的普适性, 从理论层面阐释了该现象的作用机制. 
    具体而言, 在对抗训练的最终阶段, 神经网络输出层特征的单纯形等角紧框架 (Simplex ETF) 几何尺寸较标准训练显著缩小, 且压缩程度随对抗扰动半径的增大而递增. 
    本文通过大量实验验证了该现象在多种模型架构, 数据集和对抗训练算法中的普适性, 并从特征学习动力学的角度提供了严格的理论解释, 证明了单纯形压缩是神经塌缩 (Neural Collapse) 现象在对抗训练场景下的自然推广与延伸。

    \item \textbf{建立了重构式异常检测（RBAD）的特征学习理论分析框架}: 
    本文针对仅使用正常数据训练的重构式异常检测方法, 构建了基于特征学习理论的分析框架. 创新性地提出了 ``锥集直觉'' (Cone Set Intuition) 来刻画卷积神经网络在训练过程中如何从正常数据中提取特征, 并理论证明了特征提取的动力学过程. 
    进一步地, 本文严格证明了在正常数据上训练的卷积自编码器对异常样本的重构误差显著大于正常样本, 从理论上解释了重构式异常检测方法的有效性机制. 
    该理论框架为异常检测领域提供了新的理论视角, 并在合成和真实数据集上得到了实验验证。
\end{enumerate}
